\SourceLeaf{426}%
the surface, the value of $V$ will change sign, so that as long as continuity is maintained there are positive values on one side and negative [values] on the other.

Let the point $L$ on the auxiliary sphere denote the direction of the normal to the surface toward the side we regard as upper, and let
\[
\cos L(1)=X,\qquad\cos L(2)=Y,\qquad\cos L(3)=Z.
\]
Further, let $ds$ be an infinitesimal line on the surface and, with its direction represented by the point $\lambda$ on the auxiliary sphere, let
\[
\cos\lambda(1)=\xi,\qquad\cos\lambda(2)=\eta,\qquad\cos\lambda(3)=\zeta.
\]
We then have
\[
dx=\xi ds,\qquad dy=\eta ds,\qquad dz=\zeta ds,
\]
hence
\[
P\xi+Q\eta+R\zeta=0
\]
and since $\lambda L=90^{0}$ must hold, also
\[
X\xi+Y\eta+Z\zeta=0.
\]
Since $P,Q,R,X,Y,Z$ depend only on the point of the surface from which the element is taken, and these equations hold for every direction of the element in the surface, it is easy to see that $P,Q,R$ must be proportional to $X,Y,Z$, hence
\[
P=X\mu,\qquad Q=Y\mu,\qquad R=Z\mu,
\]
Thus, since
\[
XX+YY+ZZ=1
\]
we obtain
\[
\mu=PX+QY+RZ
\]
and
\[
\mu\mu=PP+QQ+RR
\]
or
\[
\mu=\pm\BareRoot{(PP+QQ+RR)}.
\]

If we move away from the surface along the normal by\space%
